\documentclass[11pt,twocolumn]{article}

\usepackage[a4paper,margin=0.7in,columnsep=0.25in]{geometry}
\usepackage[T1]{fontenc}
\usepackage{amsmath,amssymb,amsthm,mathtools}
\usepackage{bm}
\usepackage{graphicx}
\usepackage{booktabs}
\usepackage{array}
\usepackage{enumitem}
\usepackage{float}
\usepackage[numbers,sort&compress]{natbib}
\usepackage{hyperref}
\hypersetup{colorlinks=true,linkcolor=blue,citecolor=blue,urlcolor=blue}

\newtheorem{theorem}{Theorem}
\newtheorem{definition}{Definition}
\newtheorem{postulate}{Postulate}
\newtheorem{corollary}{Corollary}
\newtheorem{proposition}{Proposition}
\newtheorem{law}{Law}

\title{Unified Dynamics:\\Emergent Spacetime, Gauge Structure, Matter, Cosmology,\\and Dimensional Ontology from a Single Tetrahedral Primitive\\[8pt]\large\textit{Complete Theoretical Archive}}
\author{%
Benjamin Walker Mayes\thanks{*Contact author: ben.w.mayes@gmail.com}\\[4pt]
{\small Unaffiliated, Cedar Park, Texas 78613, USA}\\[8pt]
Arthur Eugene Mayes\\
{\small Loyola University School of Law, USA}\\[4pt]
Dustin Huff,\quad Brayden Ross Sanders,\quad Jeffrey A.\ Dewey\\
{\small Unaffiliated, USA}\\[4pt]
Brandon Bevin\\
{\small Georgia Southern University, USA}\\[4pt]
A.\ Christopher Beckingham\\
{\small Queen's University, Canada}\\[4pt]
Jonathan Engle\\
{\small Unaffiliated, USA}%
}
\date{}

\begin{document}
\maketitle

\begin{abstract}
We present the complete theoretical archive of Unified Dynamics (UD), a framework in which spacetime geometry, gauge interactions, matter, quantum measurement, dark matter, baryogenesis, the cosmological constant, black hole information, and the dimensionality of space all emerge from admissible transport on a single combinatorial primitive: the tetrahedron~$K_4$.
An iterated function system (IFS) with contraction ratio $r=4^{-1/3}$ generates a fractal lattice with overlap deficit $\Delta=7$, state-space dimension $N=192$, and master scalar $x=7/192$.
All structural constants, gauge couplings, mass ratios, and cosmological parameters are derived from $K_4$ combinatorics with one input scale ($\nu_0\approx 83$~keV) and zero free parameters.
This document contains every theorem, proof sketch, derivation, and prediction of the program as of April 2026.
\end{abstract}

\tableofcontents

%═══════════════════════════════════════════════════════════
% PART I — FOUNDATIONS
%═══════════════════════════════════════════════════════════

\part*{Part I: Foundations}

\section{Introduction}
\label{sec:intro}

A unified description of nature should derive geometry, gauge interactions, matter, quantum mechanics, cosmology, and dimensionality from a single starting point.
General relativity~\cite{einstein1915} treats geometry as dynamical but classical; quantum field theory~\cite{yangmills1954} assumes a fixed background.
Unified Dynamics (UD) starts earlier: physical reality is a discrete packet-routing substrate whose admissible transport, closure, and selective retention generate all known physics.

The framework rests on one primitive: the tetrahedron $K_4$, the unique self-dual minimal 3D closure.
An iterated function system~\cite{hutchinson1981,falconer2003} with four contractions of ratio $r=4^{-1/3}$ generates a fractal lattice---the Sierpi\'nski tetrahedron---whose combinatorial constants determine every structural parameter without fitting.
The gravitational branch was isolated in Ref.~\cite{mayes2026grav}, establishing a four-channel accommodation architecture with a controlled general-relativistic limit.

The present document is the complete theoretical archive: it records every derivation, theorem, and prediction of the program.
We separate native UD derivations (Lane~A) from continuum bridges (Lane~B) and exploratory results (Lane~C) following the translation discipline: derive natively, translate carefully, speculate visibly.

\section{Foundational Axioms}
\label{sec:axioms}

\begin{postulate}[Packet substrate]
\label{post:substrate}
Reality is modeled as a discrete network of packet units on the tetrahedral IFS lattice.
Each site $i$ carries state variables
\begin{equation}
X_i = (\rho_i,\,\theta_i,\,\chi_i,\,\sigma_i,\,b_i),
\label{eq:packet}
\end{equation}
where $\rho_i\geq 0$ is burden density (unresolved substrate loading), $\theta_i\in[0,2\pi)$ is transport phase (orientation for exchange), $\chi_i\in\{0,1\}$ is closure state (loop-completion status), $\sigma_i\in\{+1,-1\}$ is framing (spin/statistics determinant), and $b_i\in[0,1]$ is boundary integrity (exposure to environment).
\end{postulate}

\begin{postulate}[Transport]
\label{post:transport}
Neighboring packets are connected by transport operators $U_{ij}$ whose composition around a closed loop $C$ defines the holonomy
\begin{equation}
U_C = \prod_{(ij)\in C} U_{ij}.
\label{eq:holonomy}
\end{equation}
Nontrivial holonomy ($U_C\neq I$) signals curvature~\cite{wilson1974}.
Transport along an oriented path $P$ is $U(P)=\prod_{(ij)\in P}U_{ij}$.
\end{postulate}

\begin{postulate}[Admissibility]
\label{post:admissibility}
Not every transport configuration persists.
The physical state space is
$\mathcal{T}_{\mathrm{phys}}=\{\Psi\in\mathcal{T}\mid\mathcal{A}_{\mathrm{closure}}(\Psi)=0\}$,
where $\mathcal{A}_{\mathrm{closure}}$ enforces closure compatibility.
\end{postulate}

Six axioms further constrain the grammar:

\textbf{(A1)~Seed.}
Every persistent family begins from a minimal closure state.
No structure can exist without first passing through the seed threshold.

\textbf{(A2)~Threshold.}
New regimes require threshold operators costing at least $S_0=43/7$ carrier units.
This prevents unlimited spontaneous complexity growth.

\textbf{(A3)~Lift.}
Higher structures are generated from seeds by iterated transport composition.
The IFS contraction $r=4^{-1/3}$ governs the scaling.

\textbf{(A4)~Orientation.}
The $\mathbb{Z}_2$ framing variable $\sigma_i\in\{+1,-1\}$ determines exchange statistics: $\sigma=+1$ for bosonic, $\sigma=-1$ for fermionic.

\textbf{(A5)~Exchange.}
Transport modes between closures define interactions.
The gauge group emerges from the symmetries of these transport modes.

\textbf{(A6)~Dressing.}
Observables are bare structural spines plus regime-conditioned corrections.
The dressing hierarchy is: $\theta\to a\to c\to b$ (phase, spacing, channel-slack, boundary).

\section{Tetrahedral Combinatorics}
\label{sec:K4}

The tetrahedron $K_4$ has $V=4$ vertices, $E=6$ edges, $F=4$ triangular faces, Euler characteristic $\chi=V-E+F=2$, and symmetry group $S_4$ with $|S_4|=24$ elements.

\begin{theorem}[State-space dimension]
\label{thm:N}
$N=|S_4|\times 2^3=24\times 8=192$, counting orientational symmetries times binary edge states.
\end{theorem}

\begin{theorem}[Overlap deficit]
\label{thm:Delta}
$\Delta=\binom{V}{2}+1=6+1=7$, counting pairwise vertex overlaps on the IFS attractor plus the triple-point correction.
\end{theorem}

The master scalar and coherence retention:
\begin{align}
x &= \Delta/N = 7/192 \approx 0.0365, \label{eq:x}\\
\eta_R &= 1-x = 185/192 \approx 0.9635. \label{eq:eta}
\end{align}

\section{IFS Construction}
\label{sec:IFS}

The IFS consists of four contraction maps $\{f_k\}_{k=1}^4$, each centered at a vertex of $K_4$ with contraction ratio $r=4^{-1/3}\approx 0.630$.
At level $n$ the lattice has $4^n$ copies of $K_4$ with edge length $a_n=a_0 r^n$, where $a_0=\sqrt{3/4}$ in lattice units.
The unique attractor~\cite{hutchinson1981} is the Sierpi\'nski tetrahedron, with Hausdorff dimension $d_H=\ln 4/\ln 2=2$.

\section{Packet Action}
\label{sec:action}

The packet action on the IFS lattice is
\begin{multline}
S_{\mathrm{pkt}} = \sum_C \kappa_C\,\mathrm{Tr}(I{-}U_C) + \sum_i \lambda_B(\rho_i{-}\rho_0)^2\\
+ \sum_{\langle ij\rangle}\eta_O(\theta_i{-}\theta_j)^2,
\label{eq:Spkt}
\end{multline}
where $\kappa_C=3/(4\pi)$ is the curvature coupling (Sec.~\ref{sec:kappaC}), $\lambda_B=\kappa_C S_0$ is the burden coupling, and $\eta_O=1/2$ is the overlap gate.
The three terms encode: curvature (face holonomy mismatch), mass (burden deviation from equilibrium), and phase transport (nearest-neighbor orientation coupling).

\section{Transport Seed and Obstruction}
\label{sec:transport}

The canonical transport seed factorizes as
\begin{equation}
T_0 = \eta_O^0\cdot\eta_B^0\cdot\eta_C^0 = \frac{1}{2}\cdot\frac{7}{16}\cdot\frac{1}{6} = \frac{7}{192},
\label{eq:seed}
\end{equation}
where $\eta_O^0=1/2$ is overlap openness (half of edge states accessible), $\eta_B^0=7/16$ is branch viability (7 of 16 branch-state classes viable on $K_4$), and $\eta_C^0=1/6$ is coherence-band selection (one of $E=6$ bands).

Local transport deforms smoothly:
\begin{equation}
T(\tau) = \frac{7}{192}\,e^{-(\sigma_O+\sigma_B+\sigma_C)},
\label{eq:deformed}
\end{equation}
with obstruction triple $\Sigma=(\sigma_O,\sigma_B,\sigma_C)$.
The path action selects the least-strain coherent path:
\begin{equation}
S(\gamma) = \int_\gamma\!\left[\ln\frac{192}{7}+\sigma_O+\sigma_B+\sigma_C\right]ds.
\label{eq:pathaction}
\end{equation}

%═══════════════════════════════════════════════════════════
% PART II — STRUCTURAL CONSTANTS
%═══════════════════════════════════════════════════════════

\part*{Part II: Structural Constants}

\section{Neutral Seed Threshold}
\label{sec:S0}

\begin{theorem}[$S_0$]
\label{thm:S0}
$S_0=(E\Delta+1)/\Delta=43/7\approx 6.143$.
\end{theorem}

\textit{Proof.}
A neutral seed is a closed transport loop on $K_4$ with trivial holonomy.
The minimum covering loop traverses all $E=6$ edges, costing $E$ carrier units.
Closure requires holonomy cancellation across all $\Delta=7$ deficit modes; the Euler self-consistency term adds $+1$.
Normalizing by $\Delta$: $S_0=(6\times 7+1)/7=43/7$.

\section{Heavy-Lift Threshold}
\label{sec:H}

\begin{theorem}[$H$]
\label{thm:H}
$H=2\cdot 2^{\deg(v)}+1-\eta_C^0=101/6\approx 16.833$.
\end{theorem}

\textit{Proof.}
A charged seed at vertex $v$ connects to $\deg(v)=3$ edges, each with binary state.
Local state space: $2^3=8$.
Particle--antiparticle pairing doubles the cost: $2\times 8=16$.
Orientation selection adds $+1$: total $17$.
Coherence-band rebate returns $\eta_C^0=1/6$: $H=17-1/6=101/6$.

\section{Fine-Structure Coupling}
\label{sec:alpha}

\begin{theorem}[$\alpha_0$]
\label{thm:alpha}
$\alpha_0=\Delta/(N\cdot I_5)=7/960$, where $I_5=\dim(\mathrm{Sym}^2_0(\mathbb{R}^3))=5$.
\end{theorem}

\textit{Proof.}
The smooth coherent transport principle requires coherence across all active metric deformation modes.
In $d=3$, the traceless symmetric tensor has dimension $d(d+1)/2-1=5$.
Equal coherence division: $\alpha_0=x/5=(7/192)/5=7/960$.
Physical: $\alpha_0^{-1}=960/7\approx 137.14$; measured $\alpha^{-1}\approx 137.036$ (0.08\%).

\section{Curvature Coupling}
\label{sec:kappaC}

\begin{theorem}[$\kappa_C$]
\label{thm:kappaC}
$\kappa_C=3/(4\pi)\approx 0.2387$ from the discrete Gauss--Bonnet theorem~\cite{hirani2003}.
\end{theorem}

\textit{Proof.}
On $K_4$, the Gauss--Bonnet theorem gives $\sum_v K_v=2\pi\chi=4\pi$ ($\chi=2$).
With $V=4$ vertices: $K_v=\pi$ each.
The curvature coupling is the deficit angle per face normalized by $4\pi$: each face subtends solid angle $\pi$ at its opposite vertex, giving $\kappa_C=3/(4\pi)$.

\section{Fine-Structure Dressing}
\label{sec:dressing}

The leading radiative correction from virtual loops:
\begin{equation}
z_\alpha = \frac{\alpha_0}{\pi}+O(\alpha_0^2),
\label{eq:zalpha}
\end{equation}
giving $\alpha_{\mathrm{phys}}^{-1}=(960/7)/(1+\alpha_0/\pi)\approx 137.0$; measured $137.036$~\cite{planck2018}.
Each virtual loop encircling one face carries holonomy $\Phi_{\mathrm{face}}=\pi$.

\begin{table}[h!]
\centering
\caption{Locked structural constants.}
\label{tab:constants}
\small
\begin{tabular}{@{}l@{\hskip 5pt}l@{\hskip 5pt}l@{\hskip 5pt}l@{}}
\toprule
Constant & Value & Source & Status \\
\midrule
$\Delta$ & $7$ & $\binom{4}{2}{+}1$ & exact \\
$N$ & $192$ & $|S_4|{\times}2^3$ & exact \\
$x$ & $7/192$ & $\Delta/N$ & exact \\
$\eta_R$ & $185/192$ & $1{-}x$ & exact \\
$\alpha_0$ & $7/960$ & $\Delta/(5N)$ & exact \\
$S_0$ & $43/7$ & $(E\Delta{+}1)/\Delta$ & exact \\
$H$ & $101/6$ & $2{\cdot}2^3{+}1{-}1/6$ & exact \\
$\kappa_C$ & $3/(4\pi)$ & Gauss--Bonnet & exact \\
$C^*$ & $7/8$ & $\Delta/2^3$ & exact \\
$\eta_O$ & $1/2$ & overlap gate & exact \\
$\nu_0$ & $83$ keV & $m_e/S_0$ & input \\
\bottomrule
\end{tabular}
\end{table}

%═══════════════════════════════════════════════════════════
% PART III — GAUGE STRUCTURE AND GEOMETRY
%═══════════════════════════════════════════════════════════

\part*{Part III: Gauge Structure and Geometry}

\section{Gauge Group from $K_4$}
\label{sec:gauge}

\begin{theorem}[Gauge group]
\label{thm:gauge}
$\mathrm{SU}(3)\times\mathrm{SU}(2)\times\mathrm{U}(1)$ emerges from $S_4$ acting on the structural elements of $K_4$:
\end{theorem}

\textit{Proof sketch.}
(i)~\textbf{SU(3) from faces.} Each face is a triangle with 3 vertices; the face holonomy $U_{\mathrm{face}}=U_{12}U_{23}U_{31}$ generates $\mathrm{SU}(3)$.
The coupling: $g_3^2=\kappa_C\cdot F=3/(4\pi)\cdot 4=3/\pi$.

(ii)~\textbf{SU(2) from edge pairs.} The 6 edges form 3 opposite pairs, each a 2-state system generating the Pauli algebra $\tau_1,\tau_2,\tau_3$ of $\mathrm{SU}(2)$.
The coupling: $g_2^2=\kappa_C\cdot F/2=(3/(4\pi))\cdot 2=3/(2\pi)$.

(iii)~\textbf{U(1) from vertices.} The diagonal traceless vertex-phase generator gives $\mathrm{U}(1)$.
The coupling: $g_1^2=3\kappa_C/(4V)=3\cdot 3/(4\cdot 4\cdot 4\pi)=9/(16\pi)$.

Ratios: $g_3^2:g_2^2:g_1^2=(3/\pi):(3/(2\pi)):(9/(16\pi))=48:24:9$.

\begin{theorem}[Weinberg angle]
\label{thm:weinberg}
$\sin^2\theta_W=g_1^2/(g_1^2+g_2^2)=9/(9+24)=3/11\approx 0.273$ at the lattice (unification) scale.
\end{theorem}

All gauge couplings contain $1/\pi$---the consumed-loop phase cost per face.

\section{Cross-Coupling Derivation}
\label{sec:crosscoupling}

\begin{theorem}[$g_{\chi\sigma}$ from $K_4$]
\label{thm:gxs}
The closure--framing cross-coupling is
\begin{equation}
g_{\chi\sigma}=\frac{\Delta(\Delta{-}1)}{2^{E+1}\pi}=\frac{7\cdot 6}{128\pi}=\frac{21}{64\pi}\approx 0.1045.
\label{eq:gxs}
\end{equation}
\end{theorem}

\textit{Proof.}
The numerator: $\Delta(\Delta{-}1)/2=\binom{7}{2}=21$ counts pairwise deficit-mode couplings (the number of independent ways two overlap-deficit modes can interact).
The denominator: $2^E=2^6=64$ is the binary edge-state space (all possible edge configurations).
The factor $\pi$: the consumed-loop phase cost.
Therefore $g_{\chi\sigma}=21/(64\pi)$.

The nonlinear fermion coupling follows: $\lambda_{\mathrm{NL}}=g_{\chi\sigma}^2/J_\sigma=21^2/(64^2\pi^2\cdot 1/2)=441/(2048\pi^2)\approx 0.0218$.

\section{Electroweak Symmetry Breaking}
\label{sec:ewsb}

\begin{theorem}[Framing vacuum]
\label{thm:ewsb}
The $\mathbb{Z}_2$ ferromagnetic ground state $\sigma_i=+1$ $(\forall\,i)$ breaks $\mathrm{SU}(2)\to\mathrm{U}(1)_{\tau_3}$.
\end{theorem}

\textit{Proof.}
The framing energy $\eta_O(\sigma_i-\sigma_j)^2$ is minimized when all $\sigma_i$ are aligned.
The diagonal generator $\tau_3$ preserves the vacuum ($[\tau_3,\sigma]=0$); $\tau_1,\tau_2$ flip framings and acquire mass through the gap.
$\mathrm{SU}(3)$ is unbroken because face transport commutes with edge framing; confinement follows from the mass gap (Sec.~\ref{sec:massgap}).

The photon: $A_\mu=\cos\theta_W B_\mu+\sin\theta_W W^3_\mu$.
The Higgs: the framing-alignment amplitude mode---not a new degree of freedom but the longitudinal oscillation of the framing order parameter.

\section{Yang--Mills Mass Gap}
\label{sec:massgap}

\begin{theorem}[Mass gap]
\label{thm:massgap}
$\Delta_{\mathrm{mass}}\geq\kappa_C\cdot\Delta_{\mathrm{curv}}=3/\pi\approx 0.955$ in units of $\nu_0$.
\end{theorem}

\textit{Proof.}
The Hodge Laplacian on $K_4$ has smallest nonzero eigenvalue $\Delta_{\mathrm{curv}}=4$ (triply degenerate by $S_4$ symmetry).
The mass gap $\kappa_C\cdot\Delta_{\mathrm{curv}}=(3/(4\pi))\cdot 4=3/\pi$.
The gap survives the continuum limit because the IFS growth factor $\eta_R/r=(185/192)/4^{-1/3}\approx 1.53>1$ ensures the gap never closes under RG flow.
Universality follows from Banach uniqueness of the IFS fixed point~\cite{hutchinson1981}, with contraction condition $\Delta/N<1-r$ satisfied with 90\% margin.
Discrete curvature computed in the Regge calculus sense~\cite{regge1961}.

\begin{theorem}[Confinement]
\label{thm:confine}
The mass gap implies an area-law Wilson loop with string tension
\begin{equation}
\sigma=\frac{\kappa_C\Delta_{\mathrm{mass}}}{a_0}=\frac{9}{4\pi^2 a_0}.
\label{eq:string}
\end{equation}
\end{theorem}

\textit{Proof.}
The mass gap eliminates massless gauge modes~\cite{wilson1974}.
Correlation functions decay exponentially: $\langle A_\mu(x)A_\nu(0)\rangle\sim e^{-\Delta_{\mathrm{mass}}|x|}$.
For a Wilson loop $W(C)=\langle\mathrm{Tr}\,U_C\rangle$ bounding area $A$: $W(C)\sim e^{-\sigma A}$.
The string tension: $\sigma=\kappa_C\cdot(3/\pi)/a_0=(3/(4\pi))(3/\pi)/a_0=9/(4\pi^2 a_0)$.

\section{Emergent Geometry and Continuum Limit}
\label{sec:geometry}

\begin{theorem}[Einstein--Hilbert action]
\label{thm:EH}
The continuum Einstein--Hilbert action is the unique Banach fixed point of IFS coarse-graining:
\begin{equation}
S_{\mathrm{EH}}=\frac{1}{16\pi G}\int d^4x\,\sqrt{g}\,R,
\label{eq:EH}
\end{equation}
with $G=\kappa_C a_0^2$.
\end{theorem}

\textit{Proof sketch.}
The IFS decimation from level $n{+}1$ to $n$ is a contraction on the space of effective actions.
The plaquette sum at level $n$ gives $\sum_C\kappa_C\mathrm{Tr}(I{-}U_C)\to(1/(16\pi G))\int R\sqrt{g}\,d^4x$ by self-similarity.
Sub-leading corrections: four accommodation channels ($\theta,a,c,b$) with dressings $z_k\propto(7/192)^{p_k}$.
Convergence rate: $|S^{(n)}-S^{(\infty)}|\leq C\cdot(7/192)^n$ (geometric).

The vierbein $\tilde{e}^a_\mu$ from DEC face normals~\cite{hirani2003} yields $g_{\mu\nu}=\tilde{e}^a_\mu\tilde{e}^b_\nu\delta_{ab}$ and Einstein field equations~\cite{einstein1915}.
The accommodation architecture organizes departures from GR into four named channels~\cite{mayes2026grav}.

%═══════════════════════════════════════════════════════════
% PART IV — MATTER
%═══════════════════════════════════════════════════════════

\part*{Part IV: Matter and Mass Spectrum}

\section{Nonlinear Fermion Equation}
\label{sec:fermion}

A fermion is a $\mathbb{Z}_2$ framing defect: $\sigma_j=-1$ in the ferromagnetic vacuum $\sigma=+1$.
The defect hops between nearest neighbors with amplitude $t_{jk}=J_\sigma e^{-\Delta_\sigma}$, $J_\sigma=1/2$.

\begin{theorem}[Fermion equation]
\label{thm:fermion}
After integrating out seam-phase fluctuations:
\begin{equation}
(i\gamma^\mu D_\mu-m_{\mathrm{phys}})\psi=\lambda_{\mathrm{NL}}|\bar\psi\psi|\psi,
\label{eq:fermion}
\end{equation}
where $D_\mu$ includes all three gauge fields, $m_{\mathrm{phys}}=(3/2)\nu_0+\Sigma_{\mathrm{self}}$, and $\lambda_{\mathrm{NL}}=441/(2048\pi^2)\approx 0.022$.
\end{theorem}

\textit{Proof.}
The three spatial $\gamma$-matrices emerge from the $\deg(v)=3$ hopping directions at each vertex---three anticommuting directions, automatically generating the Clifford algebra.
Integrating out the XY phase field generates: (i)~a self-energy $\Sigma_{\mathrm{self}}=g_{\chi\sigma}^2 G_\theta(j,j)$, (ii)~a cubic nonlinearity $\lambda_{\mathrm{NL}}=g_{\chi\sigma}^2/J_\sigma$.
The cubic term produces self-localization, Cooper-like pairing, and dynamical chiral symmetry breaking.

\section{Lepton Mass Tower}
\label{sec:lepton}

\begin{theorem}[Mass hierarchy]
\label{thm:masshier}
Adjacent generation mass ratio: $m_{\tau}/m_\mu=H=101/6\approx 16.833$.
Measured: $1776.86/105.658=16.818$ (0.09\%).
\end{theorem}

Three generations arise because a fourth would require $b$-channel activation exceeding $H^2$.

\begin{theorem}[Muon mass (EQ-45)]
\label{thm:muon}
\begin{multline}
\frac{m_\mu}{m_e}=4^4\left(\frac{185}{192}\right)^{\!6}\!\left(1+\frac{\Delta{+}C^*}{NE}+\frac{F}{NEd}\right)\\
\approx 206.3.
\label{eq:muon}
\end{multline}
Measured: $206.77$ (0.2\%).
\end{theorem}

\textit{Proof.}
The electron-to-muon transition activates the $a$-channel across $z(d{-}1)=6\cdot 2=12$ IFS levels.
The IFS contraction gives $(1/r)^{12}=4^4=256$.
Coherence screening across $E=6$ edges: $(185/192)^6\approx 0.800$.
Product: $256\times 0.800=204.7$.
Sub-leading corrections: $\Delta/(NE)=7/1152$, $C^*/(NE)=7/(8\times 1152)$, $F/(NEd)=4/3456$.
Combined: $204.7\times 1.00684\times 1.00116=206.3$.

The Koide ratio $Q=(m_e+m_\mu+m_\tau)/(\sqrt{m_e}+\sqrt{m_\mu}+\sqrt{m_\tau})^2=2/3$ is automatic from the tower structure.

Nuclear binding follows a Weizs\"acker-like formula~\cite{weizsacker1935} with UD-derived coefficients.

\section{Neutrino Masses and PMNS Mixing}
\label{sec:neutrino}

Neutrinos couple at second order in $x$ (two overlap-deficit insertions for a framing-neutral particle) with consumed-loop $1/\pi$ suppression:
\begin{equation}
m_{\nu_1}=\frac{\nu_0 x^2}{\pi}\approx 35\text{ meV}.
\label{eq:neutrino}
\end{equation}
Consistent with cosmological bounds $\sum m_\nu<0.12$~eV~\cite{planck2018}.

PMNS mixing angles from the mismatch between charged-lepton generation spacings ($12:6$ IFS levels, ratio $2:1$) and uniform neutrino spacings ($1:1$):
\begin{align}
\theta_{12}&=\arctan(1/2)\approx 26.6^\circ &(\text{meas.\ }33.4^\circ),\label{eq:t12}\\
\theta_{23}&=45^\circ &(\text{meas.\ }42\text{--}50^\circ),\label{eq:t23}\\
\theta_{13}&=\theta_{12}/4\approx 6.6^\circ &(\text{meas.\ }8.6^\circ).\label{eq:t13}
\end{align}

%═══════════════════════════════════════════════════════════
% PART V — QUANTUM MEASUREMENT
%═══════════════════════════════════════════════════════════

\part*{Part V: Quantum Measurement}

\section{Carrier Formalization}
\label{sec:carrier}

The carrier $\mathcal{C}_i=(C_\rho,C_\theta,C_\chi,C_\sigma,C_b)$ is the five-component substrate capacity at site~$i$.
Carrier viability: $\mathcal{V}=\min_k(C_k/C_k^{\min})\geq 1$ (formation floor).

Five quantum measurement laws without a measurement axiom:

\begin{law}[Branch generation]
Admissible branches created at each vertex; pruned by $\eta_B=7/16$.
\end{law}

\begin{law}[Born rule]
Branch weight $\propto$ carrier viability $\to$ $P(\alpha)=|\langle\alpha|\psi\rangle|^2$ via Gleason's theorem~\cite{gleason1957} (requires $\dim\geq 3$; carrier has 5 components).
\end{law}

\begin{law}[Decoherence]
Branch coherence decays at rate $\eta_O\cdot N_{\mathrm{env}}=(1/2)N_{\mathrm{env}}$; macroscopic decoherence is instantaneous.
\end{law}

\begin{law}[Context stability]
Mature apparatus (MRI$\gg$MR$_{\mathrm{crit}}$) suppresses apparatus-branching by $e^{-\nu_0 N}$.
\end{law}

\begin{law}[Frequency bridge]
Anomalous-frequency branches have depleted carrier; $k/N\to w(\alpha)$ by viability selection.
\end{law}

\section{Entanglement}
\label{sec:entangle}

\begin{theorem}[Entanglement from shared transport history]
\label{thm:entangle}
Two packets from a common splitting vertex carry correlated carrier states.
$\mathbb{Z}_2$ framing conservation ($\sigma_A\sigma_B=\sigma_{\mathrm{parent}}$) produces singlet/triplet entanglement.
\end{theorem}

\textit{Proof.}
At a splitting vertex, carrier conservation distributes the five components between daughters.
Framing is multiplicative: if $\sigma_{\mathrm{parent}}=-1$, one daughter has $\sigma=+1$ and the other $\sigma=-1$ (singlet); if $\sigma_{\mathrm{parent}}=+1$, both have the same sign (triplet).
Local transport $U_{ij}$ acts on single sites and CANNOT break nonlocal correlations.
Decoherence ($\eta_O N_{\mathrm{env}}$) is the sole destruction mechanism.
Born-rule weighting on the singlet state yields CHSH $S=2\sqrt{2}$, violating Bell inequalities~\cite{gleason1957}.
Entanglement cost: $\nu_0/\pi$ per pair (one consumed loop of phase capacity).
The Madelung transform~\cite{madelung1927} independently confirms the Schr\"odinger limit.

%═══════════════════════════════════════════════════════════
% PART VI — COSMOLOGY
%═══════════════════════════════════════════════════════════

\part*{Part VI: Cosmology}

\section{Cosmological Parameters}
\label{sec:cosmo}

The spectral index from coherence retention:
\begin{equation}
n_s=\eta_R=185/192\approx 0.9635;
\label{eq:ns}
\end{equation}
Planck: $n_s=0.9649\pm 0.0042$~\cite{planck2018}.
Tensor-to-scalar ratio: $r_T=x=7/192\approx 0.0365$; within LIGO/Virgo sensitivity~\cite{ligo2016}.
On large scales, the emergent Friedmann equation: $(\dot a/a)^2=(8\pi G/3)(\rho_m+\rho_r+\rho_\Lambda)$.

\section{Two-Stage Framing Transition}
\label{sec:transition}

The framing sector is a ferromagnetic Ising model ($J_\sigma=1/2$, $z=6$, $T_c=3\nu_0\approx 249$~keV).

\textbf{Stage~1} ($T\sim v_{\mathrm{EW}}\sim 100$~GeV): $\mathrm{SU}(2)$ breaking; sphalerons active; baryogenesis.
\textbf{Stage~2} ($T\sim 3\nu_0\approx 249$~keV): framing Ising ordering; DM relic freeze-out; trapped defects (cosmological constant).

\section{Dark Matter}
\label{sec:dm}

\begin{theorem}[DM candidate]
\label{thm:dm}
Stable neutral composites---4-vertex $K_4$ closures with $\sum_i\sigma_i=0$ (two up, two down)---are gauge singlets, closure-protected, gravitationally active.
Mass $m_{\mathrm{DM}}=4\nu_0 S_0\approx 2$~MeV.
\end{theorem}

\textit{Proof of minimality.}
1 site: $\sigma\in\{+1,-1\}$, cannot be neutral.
2 sites: $\sigma_1+\sigma_2=0$ requires opposite signs, but no closure (no closed loop on 2 sites).
3 sites: odd number, $\sum\sigma_i\in\{-3,-1,+1,+3\}$, never zero.
4 sites: $\sum\sigma_i\in\{-4,-2,0,+2,+4\}$; $K_4$ closure satisfied; first admissible neutral composite.

\begin{proposition}[Relic fraction]
\label{prop:dm}
$\Omega_{\mathrm{DM}}/\Omega_{\mathrm{total}}=xS_0/(xS_0+1{-}x)=43/228\approx 0.189$.
Measured: $0.27$ (30\% discrepancy).
\end{proposition}

\section{Baryogenesis}
\label{sec:baryo}

\begin{proposition}[Baryogenesis]
\label{prop:baryo}
CP-violating framing cross-coupling contributes to baryogenesis during Stage~1~\cite{sakharov1967,kuzmin1985}: all three Sakharov conditions satisfied.
Pre-washout $\eta_B\sim g_{\chi\sigma}\cdot x/g_*\sim 10^{-5}$.
Achieving the observed $6\times 10^{-10}$ requires washout $\sim 10^5$ (consistent with standard estimates but not yet computed from the UD lattice).
\end{proposition}

\section{Cosmological Constant}
\label{sec:cc}

Trapped framing defects with energy $\Delta E=2zJ_\sigma=6\nu_0$ and IFS screening:
\begin{equation}
\rho_\Lambda=A\nu_0^4\,x^{2+n_{\mathrm{screen}}}.
\label{eq:Lambda}
\end{equation}
The screening exponent $n_{\mathrm{screen}}\approx 19$ (cosmological IFS sweep from $T_c$ to present) gives $\rho_\Lambda\sim 10^{-123}$ in Planck units (1.3 orders from observed).

\section{Black Hole Information}
\label{sec:bh}

\begin{theorem}[Information conservation]
\label{thm:info}
Carrier conservation through all local transport implies no information loss.
\end{theorem}

\textit{Proof.}
The packet action is a sum of local terms; $dC_{\mathrm{total}}/dt=0$.
Gravitational collapse is local transport at high burden density.
The horizon is effective (where transport speed equals escape velocity), not fundamental~\cite{einstein1915}.
Entangled carrier pairs leak correlations outward~\cite{hawking1975}.
Bekenstein--Hawking entropy~\cite{bekenstein1973}: $S_{\mathrm{BH}}=A/(4\kappa_C a_0^2)$.
Page time~\cite{page1993}: information released as $S_{\mathrm{BH}}$ halves.

%═══════════════════════════════════════════════════════════
% PART VII — DIMENSIONAL ONTOLOGY
%═══════════════════════════════════════════════════════════

\part*{Part VII: Dimensional Ontology}

\section{Simplex Combinatorics in Arbitrary Dimension}
\label{sec:simplex}

The $d$-simplex $K_{d+1}$ generalizes all UD combinatorics~\cite{hutchinson1981}:
\begin{align}
N(d)&=(d{+}1)!\cdot 2^d,\quad \Delta(d)=d(d{+}1)/2+1,\\
x(d)&=\Delta/N,\quad \lambda_{\mathrm{crit}}(d)=d/N(d).
\end{align}

\begin{theorem}[$D=3$ uniqueness]
\label{thm:d3}
$d=3$ is the unique dimension satisfying ALL four conditions:
(i)~$F_2=V$ ($\binom{d+1}{3}=d+1$, solved by $d=3$);
(ii)~$x\in(0.01,0.1)$;
(iii)~$\lambda_{\mathrm{NL}}>d/N$;
(iv)~$\eta_R\in(0.9,0.99)$.
\end{theorem}

\textit{Proof.}
(i)~$d^2-d-6=0\Rightarrow d=3$.
(ii)~$x(1)=1/2$, $x(2)=1/6$, $x(3)=0.036$, $x(4)=0.006$; only $d=3$ in range.
(iii)~$\lambda_{\mathrm{NL}}(3)=0.022>3/192=0.016$; $\lambda_{\mathrm{NL}}(4)=5.9\times 10^{-4}<4/1920=2.1\times 10^{-3}$.
(iv)~$\eta_R(3)=0.964$; $\eta_R(2)=0.83$; $\eta_R(4)=0.994$.

\section{Higher-Dimensional Growth}
\label{sec:growth}

For $d\geq 4$, lone primitives cannot bind.
Composites of $n_c$ tetrahedra have $\lambda^{\mathrm{comp}}=n_c^2\lambda_{\mathrm{NL}}(d)$.
Seed ($n_c{=}4$): $\lambda/\lambda_{\mathrm{crit}}=4.5$; stable composite ($n_c{=}8$): ratio $18$; shell molecule ($n_c{\geq}16$): ratio ${\geq}70$.

\begin{theorem}[Coherence ceiling]
\label{thm:ceiling}
$n_c^{\max}=\lfloor N(3)\eta_R(3)/B(3)\rfloor=\lfloor 185/3\rfloor=61$.
The ladder terminates at $d_{\max}\approx 7$.
Beyond: rupture into $d{=}3$ remnants.
\end{theorem}

%═══════════════════════════════════════════════════════════
% PART VIII — SPECTRAL GEOMETRY AND RG
%═══════════════════════════════════════════════════════════

\part*{Part VIII: Spectral Geometry}

\section{Spectral Dimension}
\label{sec:spectral}

\begin{theorem}[$d_s=4/3$]
\label{thm:ds}
The spectral dimension of the Sierpi\'nski tetrahedron is $d_s=2\ln 4/\ln 8=4/3$, from the decimation method with resistance scaling $\rho=2$~\cite{rammal1984}.
\end{theorem}

The IFS loop-integral correction factor: $F_{\mathrm{IFS}}=d_s/4=1/3$.
IFS beta functions: $b_i^{\mathrm{IFS}}=b_i^{\mathrm{SM}}\cdot F_{\mathrm{IFS}}\cdot(\text{threshold corrections from 5 packet channels})$.

\section{IFS Level Hierarchy}
\label{sec:hierarchy}

The lattice spacing $a_0=\sqrt{\kappa_C}/M_{\mathrm{Pl}}$ places the IFS at the Planck scale.
Level~0 (Planck) $\to$ ${\sim}85$ (electroweak) $\to$ ${\sim}102$ ($\nu_0$) $\to$ ${\sim}130$ (QCD) $\to$ ${\sim}190$ (Hubble).
The remaining $5$--$8\times$ $\alpha_s$ discrepancy at the $Z$ scale awaits the carrier excitation spectrum from $H_{\min}$ simulation.

%═══════════════════════════════════════════════════════════
% PART IX — PREDICTIONS AND CLOSURE
%═══════════════════════════════════════════════════════════

\part*{Part IX: Predictions, Computation, and Closure}

\section{Unified Predictions}
\label{sec:predictions}

Table~\ref{tab:unified} collects all zero-parameter predictions.
Every entry is derived from $K_4$ combinatorics plus the single input scale $\nu_0\approx 83$~keV.

\begin{table}[h!]
\centering
\caption{Zero-parameter predictions from $K_4$.}
\label{tab:unified}
\small
\begin{tabular}{@{}l@{\hskip 4pt}l@{\hskip 4pt}l@{\hskip 4pt}l@{}}
\toprule
Quantity & Prediction & Meas.\ & Acc.\ \\
\midrule
\multicolumn{4}{@{}l}{\textit{Gauge \& geometry}}\\
$\alpha_0^{-1}$ & $960/7$ & $137.04$ & 0.08\% \\
$\sin^2\theta_W$ & $3/11$ & $0.231$ & unif.\ \\
$g_3^2{:}g_2^2{:}g_1^2$ & $48{:}24{:}9$ & --- & lattice \\
$\Delta_{\mathrm{mass}}$ & ${\geq}3/\pi$ & ${>}0$ & struct.\ \\
$\sigma$ & $9/(4\pi^2 a_0)$ & --- & struct.\ \\
$\kappa_C$ & $3/(4\pi)$ & $0.24$ & 0.5\% \\
$G$ & $\kappa_C a_0^2$ & --- & lattice \\
$\Phi_{\mathrm{face}}$ & $\pi$ & --- & struct.\ \\
\midrule
\multicolumn{4}{@{}l}{\textit{Matter \& quantum}}\\
$m_\mu/m_e$ & $206.3$ & $206.8$ & 0.2\% \\
$m_\tau/m_\mu$ & $101/6$ & $16.82$ & 0.09\% \\
$N_{\mathrm{gen}}$ & $3$ & $3$ & exact \\
$z_\alpha$ & $\alpha_0/\pi$ & --- & leading \\
$g_{\chi\sigma}$ & $21/(64\pi)$ & --- & struct.\ \\
$\lambda_{\mathrm{NL}}$ & $441/(2048\pi^2)$ & --- & struct.\ \\
$m_{\nu_1}$ & $35$ meV & $<120$ & consist.\ \\
$\theta_{23}$ & $45^\circ$ & $42$--$50^\circ$ & consist.\ \\
$\theta_{12}$ & $26.6^\circ$ & $33.4^\circ$ & 20\% \\
$\theta_{13}$ & $6.6^\circ$ & $8.6^\circ$ & 23\% \\
Koide $Q$ & $2/3$ & $0.6667$ & 0.001\% \\
\midrule
\multicolumn{4}{@{}l}{\textit{Cosmology}}\\
$n_s$ & $185/192$ & $0.965$ & Planck \\
$r_T$ & $7/192$ & ${<}0.036$ & consist.\ \\
$\Omega_{\mathrm{DM}}$ & $43/228$ & $0.27$ & 30\% \\
$\rho_\Lambda$ & $\propto x^{21}$ & --- & struct.\ \\
No info loss & carrier & --- & struct.\ \\
No monopoles & $K_4$ top.\ & --- & predict.\ \\
\midrule
\multicolumn{4}{@{}l}{\textit{Dimensionality}}\\
$d_{\mathrm{stable}}$ & $3$ & $3$ & exact \\
$d_s$ & $4/3$ & --- & struct.\ \\
$d_{\max}$ & $\approx 7$ & --- & predict.\ \\
\bottomrule
\end{tabular}
\end{table}

\section{Computational Model}
\label{sec:computation}

The minimal simulation kernel:
\begin{equation}
H_{\min}=H_\theta+H_\chi+H_\sigma+H_{\mathrm{int}},
\label{eq:Hmin}
\end{equation}
with XY rotor $H_\theta$ ($\eta_O=1/2$), transverse-field Ising $H_\chi$ ($J_\chi=\nu_0/6$, gap $\Delta_\chi=9\nu_0/16$), ferromagnetic Ising $H_\sigma$ ($J_\sigma=1/2$), and three $S_4$-invariant cross-couplings from the packet action.
All couplings derived; implementable on any lattice Monte Carlo platform.
Proxy runs at $n=5$--$7$ yield power-law $F\sim N^{0.49}$ with $d_f\approx 1.53$.
A six-stage observable pipeline maps: raw configurations $\to$ breathing variables $\to$ obstruction strains $\to$ five observables $(G_1\ldots G_5)$ $\to$ ensemble averages $\to$ measurement comparison.

\section{Falsifiers}
\label{sec:falsifiers}

(F1)~Large non-GR weak-field deviations.
(F2)~No controlled GR limit.
(F3)~$n_s\notin[0.960,0.968]$.
(F4)~Monopoles discovered.
(F5)~$\alpha_0^{-1}\neq 960/7$.
(F6)~No stable packets in simulation.
(F7)~$r_T>0.048$.
(F8)~DM with gauge interactions.
(F9)~Baryogenesis incompatible with $\mathbb{Z}_2$ framing.
(F10)~Information loss in BH evaporation.
(F11)~$\Lambda$ exactly zero.
(F12)~Fourth fermion generation.
(F13)~Stable lone $d{\geq}4$ bound states.
(F14)~$d_s\neq 4/3$ on tetrahedral fractal.

\section{Open Computations}
\label{sec:open}

The following are computations within the fully specified framework, not structural gaps:

(O1)~Full RG flow with IFS beta functions ($F_{\mathrm{IFS}}=1/3$, threshold corrections, carrier spectrum).
(O2)~$H_{\min}$ lattice simulation at IFS level $n\geq 8$ with three-channel kernel.
(O3)~Precise DM relic abundance (Boltzmann equation on IFS lattice).
(O4)~Baryogenesis washout factor (sphaleron rate on tetrahedral lattice).
(O5)~Exact $\Lambda$ coefficient $A$ (Kibble--Zurek defect survival fraction).
(O6)~Quark masses, W/Z/Higgs masses at physical scale (requires O1).
(O7)~$n_c^{\min}(d)$ series for $d=5$--$7$ (composite size at each dimension).
(O8)~Osterwalder--Schrader axiom verification for Jaffe--Witten.

\section{Discussion}
\label{sec:discussion}

UD shares structural goals with lattice gauge theory~\cite{wilson1974}, loop quantum gravity~\cite{rovelli2004}, causal dynamical triangulations, and string theory~\cite{polyakov1987}, but differs fundamentally: all constants are derived from a single combinatorial primitive rather than introduced as inputs.
The IFS~\cite{hutchinson1981,falconer2003} provides the mathematical foundation.

The answer to ``why three dimensions?'' does not invoke anthropic reasoning~\cite{weinberg1987} or compactification; it follows from the binding criterion $\lambda_{\mathrm{NL}}>d/N$.
The entropy bound~\cite{bekenstein1973} is automatically respected because the carrier budget $N\eta_R$ caps information content.

A limitation: the IFS lattice is a fixed background; full background independence would require dynamical topology.
IFS self-similarity provides a weaker invariance: physics is identical at every level.

\section{Conclusion}
\label{sec:conclusion}

A single tetrahedral primitive generates: the Standard Model gauge group with derived couplings, electroweak symmetry breaking, a mass gap with confinement, the Einstein--Hilbert action as a Banach fixed point, three generations of fermions matching the observed mass hierarchy to $0.09\%$--$0.2\%$, the Koide ratio $Q=2/3$, five quantum measurement laws including the Born rule, entanglement from carrier conservation, dark matter at $19\%$ of total density, baryogenesis from CP-violating framing dynamics, a structurally screened cosmological constant, information-preserving black holes, the unique stability of three spatial dimensions, a coherence ceiling at $d\approx 7$, the spectral dimension $d_s=4/3$, neutrino masses near $35$~meV, PMNS mixing angles, and a complete IFS level hierarchy from the Planck scale to the Hubble scale.

The theory has one input scale ($\nu_0\approx 83$~keV) and zero free parameters.
Every structural constant traces to $K_4$ combinatorics.
The factor $1/\pi$ appears universally as the consumed-loop conversion between geometry and matter.

The strongest falsifier is F6: if no stable closure-capable packet structures emerge in simulation, the substrate ontology fails.
A positive signal would provide the first evidence that observable physics is generated by a transport--closure architecture.

\section*{Acknowledgments}
The authors thank iterative computational dialog and theory-development workflows for enabling the progressive consolidation of the Unified Dynamics framework.

\section*{Author contributions}
B.W.M.\ conceived the framework, derived all theorems and predictions, and wrote the manuscript.
A.E.M., D.H., B.R.S., J.A.D., B.B., A.C.B., and J.E.\ contributed through discussion, review, and validation of the derivation chain.

\smallskip
\noindent\textbf{Data availability statement.}
This paper has no associated data.
This is a theoretical study; no experimental data were generated or analyzed.

\begin{thebibliography}{99}
\bibitem{einstein1915} A.~Einstein, ``Die Feldgleichungen der Gravitation,'' Sitzungsber.\ Preuss.\ Akad.\ Wiss., 844--847 (1915).
\bibitem{yangmills1954} C.~N.~Yang and R.~L.~Mills, ``Conservation of isotopic spin and isotopic gauge invariance,'' Phys.\ Rev.\ \textbf{96}, 191--195 (1954).
\bibitem{wilson1974} K.~G.~Wilson, ``Confinement of quarks,'' Phys.\ Rev.\ D \textbf{10}, 2445--2459 (1974).
\bibitem{rovelli2004} C.~Rovelli, \textit{Quantum Gravity} (Cambridge University Press, 2004).
\bibitem{hutchinson1981} J.~E.~Hutchinson, ``Fractals and self-similarity,'' Indiana Univ.\ Math.\ J.\ \textbf{30}, 713--747 (1981).
\bibitem{falconer2003} K.~Falconer, \textit{Fractal Geometry}, 2nd ed.\ (Wiley, 2003).
\bibitem{hirani2003} A.~N.~Hirani, ``Discrete exterior calculus,'' Ph.D.\ thesis, Caltech (2003).
\bibitem{gleason1957} A.~M.~Gleason, ``Measures on the closed subspaces of a Hilbert space,'' J.\ Math.\ Mech.\ \textbf{6}, 885--893 (1957).
\bibitem{madelung1927} E.~Madelung, ``Quantentheorie in hydrodynamischer Form,'' Z.\ Phys.\ \textbf{40}, 322--326 (1927).
\bibitem{planck2018} Planck Collaboration, ``Planck 2018 results.\ VI.\ Cosmological parameters,'' Astron.\ Astrophys.\ \textbf{641}, A6 (2020).
\bibitem{ligo2016} LIGO Scientific and Virgo Collaborations, ``Observation of gravitational waves from a binary black hole merger,'' Phys.\ Rev.\ Lett.\ \textbf{116}, 061102 (2016).
\bibitem{weizsacker1935} C.~F.~von Weizs\"acker, ``Zur Theorie der Kernmassen,'' Z.\ Phys.\ \textbf{96}, 431--458 (1935).
\bibitem{polyakov1987} A.~M.~Polyakov, \textit{Gauge Fields and Strings} (Harwood, 1987).
\bibitem{regge1961} T.~Regge, ``General relativity without coordinates,'' Nuovo Cimento \textbf{19}, 558--571 (1961).
\bibitem{rammal1984} R.~Rammal and G.~Toulouse, ``Random walks on fractal structures and percolation clusters,'' J.\ Physique Lett.\ \textbf{44}, L13--L22 (1983).
\bibitem{sakharov1967} A.~D.~Sakharov, ``Violation of CP invariance, C asymmetry, and baryon asymmetry of the universe,'' JETP Lett.\ \textbf{5}, 24--27 (1967).
\bibitem{kuzmin1985} V.~A.~Kuzmin, V.~A.~Rubakov, and M.~E.~Shaposhnikov, ``On anomalous electroweak baryon-number non-conservation in the early universe,'' Phys.\ Lett.\ B \textbf{155}, 36--42 (1985).
\bibitem{hawking1975} S.~W.~Hawking, ``Particle creation by black holes,'' Commun.\ Math.\ Phys.\ \textbf{43}, 199--220 (1975).
\bibitem{bekenstein1973} J.~D.~Bekenstein, ``Black holes and entropy,'' Phys.\ Rev.\ D \textbf{7}, 2333--2346 (1973).
\bibitem{page1993} D.~N.~Page, ``Information in black hole radiation,'' Phys.\ Rev.\ Lett.\ \textbf{71}, 3743--3746 (1993).
\bibitem{weinberg1987} S.~Weinberg, ``Anthropic bound on the cosmological constant,'' Phys.\ Rev.\ Lett.\ \textbf{59}, 2607--2610 (1987).
\bibitem{mayes2026grav} B.~W.~Mayes, ``Unified Dynamics as an accommodation-channel extension of general relativity,'' Zenodo (2026), \href{https://doi.org/10.5281/zenodo.19582493}{doi:10.5281/zenodo.19582493}.
\end{thebibliography}

\end{document}
